\chapter{Continuation Algorithms: Inactive Modes and Extensions}
\label{app:continuation}

This appendix chapter collects continuation branches and control-flow extensions
that remain in the repository but are not active in the committed default PMG
benchmark.

\section{Direct SSR}

The direct SSR path is still implemented through:
\begin{itemize}
\item \petscmod{continuation/direct.py},
\item \petscmod{continuation/omega.py}.
\end{itemize}

It remains the closest PETSc analogue of the simpler MATLAB direct continuation
scripts, but it is not the current default path.

\section{Limit Load}

The limit-load branch remains available through
\petscmod{continuation/limit\_load.py}. It solves a different continuation
problem with the load multiplier \(t\) instead of the SSR reduction factor
\(\lambda\). This is active code, but outside the current benchmark.

\section{Predictor Families Beyond Secant}

\petscmod{continuation/indirect.py} supports predictor families beyond the
mainline \code{secant} mode:
\begin{itemize}
\item \code{reduced\_newton\_all\_prev},
\item \code{reduced\_newton\_affine\_all\_prev},
\item \code{reduced\_newton\_window},
\item \code{reduced\_newton\_increment\_power}.
\end{itemize}

These change the quality of the initial guess, not the continuation equations.

\section{Additional Control-Flow Extensions}

The inactive feature set also includes:
\begin{itemize}
\item adaptive omega-step control,
\item secant correction modes,
\item history-box step-length caps,
\item fine-switch logic,
\item temporary first-Newton warm starts,
\item \code{streaming\_microstep} continuation mode.
\end{itemize}

These are real PETSc-era extensions, but they are not part of the committed
default 3D benchmark configuration.

\section{Why They Were Moved Out Of The Main Text}

The mainline continuation chapter now follows the exact benchmark settings:
\code{indirect + secant + legacy + classic}. The alternatives remain documented
here so the document stays useful for future feature work without obscuring what
the default benchmark actually does today.
